\chapter{Применения разработанного метода проектирования}
\begin{comment}
	\textit{Четвертая глава содержит решения конкретных задач со всеми краевыми условиями, расчет конкретного устройства, графики, зависимости, вторичные модели, результаты проверки сходимости теоретических положений с экспериментальными данными для конкретной модели и т.д.
Обсуждению и оценке результатов диссертационной работы можно посвятить отдельный параграф.
Оценка результатов работы должна быть качественной и количественной.
Сравнение с известными решениями следует проводить по всем возможным аспектам.
Следует указать на возможность обобщений, дальнейшего развития методов и идей, использования результатов диссертации в смежных областях (с соблюдением необходимой корректности).
Объем главы 20-25 страниц.}
\end{comment}

\section{Введение}
Традиционные методы проектирования осевых вентиляторов и компрессоров основаны на допущениях, приводящих к искажению характеристик вентиляторов с меридиональным ускорением потока.
Универсальные сеточные методы моделировании течения решением уравнений Навье"--~Стокса осреднённых по Рейнольдсу (RANS) не содержат этих допущений, и активно используются при анализе течения в проточной части турбомашин.
При современном развитии RANS и множеству их реализаций в виде программ на ЭВМ, методы потенциального течения не теряют свою актуальность.
Она обусловлена их низкой требовательностью к ЭВМ.
В главе приводится сравнение разработанных методов основанных на потенциальном течении с другими известными и обсуждаются возможности их практического применения.

Результаты моделирования течения сравниваются с полученными методами RANS по получаемым характеристикам вентилятора с меридиональным ускорением потока и времени затраченного на расчёт.

По соответствию распределения теоретического напора по высоте проточной части расчётному проводится сравнение разработанного квазитрёхмерного метода проектирования с основанными на пересчёте характеристик плоских решёток профилей.

В качестве применений, кроме проектирования вентиляторов с меридиональным ускорением потока, обсуждается проектирование спрямляющих аппаратов в S-образных каналах, рабочих колёс центробежных вентиляторов и осе-радиальных компрессоров.
Каждый из этих элементов лопаточных машин обладает особенностями, затрудняющими их разработку методами основанными на проектировании осевых машин.
Для примера выполняется профилирование колёс и проверочным расчётом методом RANS подтверждаются расчётные характеристики.


\section{Сравнение методов моделирования течения}\label{sec:sravenieModelirovanija}
Квазитрёхмерный метод моделирования течения позволяет с удовлетворительной точностью моделировать интегральные характеристики вентилятора с меридиональным ускорением потока и их распределение по высоте проточной части.
Влияние вязкости при этом определяется эмпирическим коэффициентом снижения циркуляции \(k_\Gamma\), учитывающим влияние пограничного слоя на течение.
Более подробную оценку этого влияния можно провести используя интегральные методы анализа пограничного слоя по распределению скорости и давления на поверхности лопатки \autocite{Schlichting1974, Truckenbrodt1952, Chen1979}.
Для этого необходимо иметь уверенность в рассчитанных распределений.

Экспериментальное определение распределений параметров по поверхности лопатки рабочего колеса "--- не тривиальная задача.
Для подтверждения расчётных характеристик полученных МДВ (раздел \ref{sec:newMDV}) проведено сравнение с характеристиками полученными более надёжным расчётным методом RANS (раздел \ref{sec:cfd}).
Так же сравнивалось время затраченное на расчёт.

Проведено сравнение методов трёх расчёта распределения давления \autocite{furashov2025}:
\begin{enumerate}
	\item методом RANS в целом межлопаточном канале;
	\item по трём осесимметричным поверхностям тока в слое переменной толщины методом RANS;
	\item по осесимметричным поверхностям тока в слое переменной толщины методом дискретных вихрей (\ref{sec:newMDV}).
\end{enumerate}

Основные параметры исследованного рабочего колеса приведены в таблице \ref{paramRotor}.
\begin{table}[htb]
	\caption{Параметры рабочего колеса.}
	\label{paramRotor}
	\centering
	\begin{tabular}{l|c|c}
		Коэффициент теоретического напора & \(\bar{H}_\text{т}\) & 0,5\\\hline
		Коэффициент осевой скорости на выходе из венца & \(\bar{c}_{2a}\) & 0,784\\\hline
		Относительный диаметр втулки перед РК &  \(\bar{d}_{1}\) & 0,4\\\hline
		Относительный диаметр втулки за РК & \(\bar{d}_{2}\) & 0,7\\\hline
		Коэффициент меридионального ускорения & \(m_{f}\) & 1,615\\\hline
		Количество лопаток & \(Z_{к}\) & 11\\\hline
		Угол атаки на средней поверхности тока, \(^{\circ}\) & \(i\) & 0\\\hline
		Густота решётки на средней линии & \(\tau\) & 1,274\\\hline
		Число Рейнольдса & \(\Re_{w_1}\) & \(4,5 \times 10^5\)\\\hline
		Частота вращения РК, об/мин & \(n\) & 1000\\\hline
		Окружная скорость, м/с & \(U_\text{к}\) & 36,65\\
	\end{tabular}
\end{table}

Формы исследуемых осесимметричных поверхностей тока у втулки, на средней линии и у периферии были определены по уравнению радиального равновесия \ref{sec:RadialEquilibrium}.
Средняя безразмерная высота лопатки \(\bar{h}=\left(r-r_\text{вт}\right)/\left(R_\text{к}-r_\text{вт}\right)\) в этих сечениях составила 0,169, 0,573, 0,901,
\begin{eqexpl}
	 \item{\(r\)}[=]\((r_{1}+r_{2})/2\) "--- средний радиус сечения лопатки, м,
	 \item{\(r_\text{вт}\)}[=]\(({r_\text{вт}}_{1}+{r_\text{вт}}_{2})/2\) "--- средний радиус втулки, м,
	 \item{\(R_\text{к}\)} радиус периферии рабочего колеса, м.
\end{eqexpl}

\begin{figure}[htb]
	\centering
	\includegraphics[]{4763.docx.tmp/word/media/image6.png}
	\caption{Расчётные поверхности тока и сечения лопатки рабочего колеса.}
	\label{fig:3section}
\end{figure}

Для RANS-расчётов строилась структурированная сетка из октаэдрических элементов.
Расчёты выполнялись при задействовании одного ядра процессора с тактовой частотой 3,2~ГГц.
Время, затраченное на расчёт целого межлопаточного канала методом RANS (1,35 млн. ячеек), составило \(\approx900\)~с.
Среднее время, затраченное на расчёт течения в одном сечении (0,1 млн. ячеек), составило 131~с.

Так как сравнивать время моделирования вязкого и невязкого газа не совсем корректно, было выполнено дополнительное моделирование течения невязкого газа сеточным методом.
Расчётная неструктурированная сетка состояла из тетраэдрических элементов размером приблизительно 250 тысяч ячеек.
Время расчёта составило 254~с.

Время расчёта распределения давления МДВ (600 вихрей) составило 0,2~с.
При этом выполнено два решения СЛАУ, позволяющих определить параметры любого режима течения как линейную комбинацию \ref{eq:lin_comb}.

На рисунке \ref{fig:p_epure} дано сравнение распределения давления на лопатках.
\begin{eqexpl}[4 mm]
	\item{\(\bar{p}\)}[=]\(\left(p-p_{1}\right)/\left(0,5\rho U_{К}^{2}\right)\)"--- приведённое статическое давление на лопатке,
	\item{\(p_{1}\)}  давление перед РК, Па,
	\item{\(p\)} давление на контуре, Па,
	\item{\(a\)} осевая координата, м.
\end{eqexpl}

\begin{figure}[htb]
	\centering
	\subcaptionbox{\(\bar{h}=0,169\)\label{fig:p_epure_a}}{
		\includegraphics[page=1]{img/ch2/p_epure}
	}
	\subcaptionbox{\(\bar{h}=0,573\)\label{fig:p_epure_b}}{
		\includegraphics[page=2]{img/ch2/p_epure}
	}
	\subcaptionbox{\(\bar{h}=0,901\)\label{fig:p_epure_c}}{
		\includegraphics[page=3]{img/ch2/p_epure}
	}
	\caption{Сравнение распределения давления на различных сечениях лопатки.}
	\label{fig:p_epure}
\end{figure}

На средней поверхности тока проточной части (рисунок \ref{fig:p_epure_b}) распределения давлений тремя способами удовлетворительно совпадают.
В периферийных и втулочных сечениях лопатки (рисунки \ref{fig:p_epure_a} и \ref{fig:p_epure_c} ) отличия \(\bar{p}\) возникают из-за влияния концевых эффектов: взаимодействий пограничных слоёв у поверхности лопатки и торцевых поверхностей, действием центробежных сил, течении газа в радиальном зазоре и перетечек в пограничном слое у поверхностей под действием градиента давлений между спинкой и корытцем соседних лопаток.
При моделировании вязкого газа в целом межлопаточном канале эти особенности концевых областей лопатки учитываются "--- слои течения оказывают взаимное влияние.

%Упрощение течения в проточной части вентилятора с меридиональным ускорением потока позволяет достаточно точно, в пределах 3\%, предсказать распределение \(\bar{H}_\text{т}\) по высоте проточной части.
%Для RANS-расчётов строилась структурированная сетка из окатаэдрических элементов.
%Расчёты выполнялись при задействовании одного ядра процессора с тактовой частотой 3,2~ГГц.
%Время, затраченное на расчёт целого межлопаточного канала методом RANS (1,4 млн. ячеек), составило 900~с.
%Среднее время, затраченное на расчёт течения в одном сечении (0,1 млн. ячеек), составило 131~с.
%Время расчёта распределения давления методом МДВ (600 вихрей) составило 0,2~с.

%Разделение проточной части коаксиальными поверхностями на независимые слои течения с использованием уравнения радиального равновесия позволило существенно (в 7 раз) снизить время моделирования течения, в сравнении с моделированием целого межлопаточного канала.
%Использование МДВ позволило сократить время моделирования в 4,5 тысяч раз в сравнении с моделированием вязкого газа и 1,26 тысяч раз невязкого.

Ценой малого времени моделирования является игнорирование большинства эффектов вязкого трёхмерного течения.
В случае, когда необходимо быстро оценивать характеристики большого количества лопаточных венцов, например при предварительном проектировании, скорость счёта определяет выбор в пользу квазитрёхмерных методов.

\section{Сравнение точности методов профилирования лопаток}\label{sec:sravenie}
Для проектирования вентилятора с меридиональным ускорением потока было использовано три подхода:
\begin{enumerate}[label=(\alph*)]
\item по результатам продувок плоских решёток профилей, игнорируя сужение проточной части \autocite{Komarov1974, Gurevich1974}.
\item по результатам продувок неподвижных решёток профилей в изолированных тонких осесимметричных слоях течения, образованных линиями тока в меридиональном сечении проточной части \autocite{Nikitin1966, Sachkova2000}.
\item так же как b, но с учётом вращения решёток профилей (раздел \ref{sec:newMDV}).
\end{enumerate}

\begin{figure} [htb]
	\centering
	\begin{subcaptionblock}{0.5\textwidth}
		%\centering
		\includegraphics[page=1]{img/ch1/2met-ab}
		\caption{С игнорированием формы проточной части}
	\end{subcaptionblock}
	\begin{subcaptionblock}{0.5\textwidth}
		%\centering
		\includegraphics[page=2]{img/ch1/2met-ab}
		\caption{С учётом перестройки потока от формы проточной части}
	\end{subcaptionblock}
	\caption{Схема расположения расчётных поверхностей тока в проточной части вентилятора.}
	\label{fig:2met-ab}
\end{figure}

Вариант (a) классический.
Преимущество подхода (b) заключено в возможности пересчёта результатов экспериментальных продувок плоских решёток профилей с изменяющейся осевой компонентой скорости с помощью конформных отображений.
Переход к неинерционной системе отсчёта при подходе (b) может привести к ошибкам проектирования ротора.
В таком случае подход c окажется предпочтителен несмотря на то, что вряд ли выйдет при этом использовать надёжные экспериментальные данные.

На одинаковые параметры (таблица \ref{tab:3RK}) были спрофилированы три рабочих колеса вентилятора с использованием перечисленных подходов.

\begin{table}[htb]
	\caption{Параметры РК}
	\label{tab:3RK}
	\begin{center}
		\begin{tabular}{l|c|c}
			Диаметр рабочего колеса, м & \(D_\text{к}\) & 0,7 \\ \hline
			Коэффициент теоретического напора& \(\bar{H}_\text{т}\) &	0,525 \\ \hline
			Коэффициент производительности&\(\phi\)	&	0,6 \\ \hline
			Коэффициент сужения площади проточной части&\(m_f\)	&	1,14 \\ \hline
			Относительный диаметр втулки на входе& \(\bar{d}_2\)	&	0,4 \\ \hline
			Относительный диаметр втулки на выходе& \(\bar{d}_2\)	&	0,52 \\ \hline
			Закон профилирования &\(c_{u}r\)&\(\const\)
		\end{tabular}
	\end{center}
\end{table}

Углы атаки изменялись от минус 2\textdegree~на периферии до 4\textdegree~на втулке.
Густоты решёток всех роторов одинаковы и определённому по \autocite{Bunimovich1967} для ступени с цилиндрической втулкой.
Углы атаки также были одинаковы, от минус 2\textdegree~на периферии до 4\textdegree~на втулке.
Параметр снижения циркуляции \(k_\Gamma\) был определён по предварительному RANS-расчёту вязкого течения с замыканием моделью турбулентностью \(k\) "--- \(\omega\) в элементарном венце.
Были уточнены значения \(k_\Gamma\) на периферийном, среднем и втулочном сечении, в остальных сечениях получены интерполяцией.
\begin{table}[htb]
\caption{Параметры решёток профилей роторов выполненых тремя методами}
\label{tab:3met}
\begin{center}
	\begin{tabular}{r|c|c|c|c|c|c|c|c}
		& & & \multicolumn{2}{c|}{а} & \multicolumn{2}{c|}{b} & \multicolumn{2}{c}{с} \\
		\hline
		\textnumero~сечения  &\(i\), \textdegree & \(\tau\) & \(\bar{f}\) & \(\theta_{\text{г}}\)& \(\bar{f}\) & \(\theta_{\text{г}}\)  & \(\bar{f}\) & \(\theta_{\text{г}}\) \\
		\hline\hline
	\hline
	1 & 1.44 & \(-2.0\) & 39.5 & 0.062 & 39.3 & 0.074 & 39.2 & 0.071 \\
	\hline
	2 & 1.46 & \(-1.2\) & 35.9 & 0.071 & 36.3 & 0.077 & 35.8 & 0.076 \\
	\hline
	3 & 1.49 & \(-0.5\) & 31.5 & 0.083 & 32.9 & 0.081 & 31.9 & 0.083 \\
	\hline
	4 & 1.52 & 0.2 & 26.3 & 0.098 & 28.9 & 0.086 & 27.3 & 0.091 \\
	\hline
	5 & 1.54 & 1.0 & 20.0 & 0.116 & 24.1 & 0.092 & 21.9 & 0.102 \\
	\hline
	6 & 1.52 & 1.7 & 12.4 & 0.139 & 18.6 & 0.100 & 15.6 & 0.116 \\
	\hline
	7 & 1.53 & 2.5 & 3.5 & 0.166 & 11.8 & 0.110 & 8.1 & 0.133 \\
	\hline
	8 & 1.59 & 3.2 & -6.4 & 0.192 & 3.2 & 0.121 & -0.9 & 0.153 \\
	\hline
	9 & 1.65 & 4.0 & -17.7 & 0.217 & -7.5 & 0.137 & -11.5 & 0.179 \\
	\hline
	\end{tabular}
\end{center}
\end{table}

Распределение циркуляции по высоте проточной части полученное в расчёте RANS на рисунке \ref{fig:3met}.
При профилировании лопаток с игнорированием формы проточной части (а) распределение \(\Ht\) занижено относительно расчётного значения по всей высоте проточной части, интегральный \(\Ht\) оказался на 3,8\% ниже расчётного.
Учёт перестройки потока из-за изменения диаметра втулки и использование конформных отображений (b) привело к ещё большему снижению \(\Ht\) до 10\% меньше расчётного.
Использование разработанного метода моделирования позволило повысить точность задания \(\Ht\) и уменьшить разницу расчётного и проверенного значений до 0,2\%.
\begin{figure}[htb]
	\centering
	\includegraphics{ch1/3metoden}
	\caption{Сравнение распределения \(\Ht\) по высоте проточной части \(\bar{h}\) рабочих колёс с лопатками профилированными тремя методами.}
	\label{fig:3met}
\end{figure}

%Результат аналогичен случаю с невязким газом.
Опираясь на результаты моделирования предварительного проекта ротора полученным любым из указанных подходов можно уточнить геометрию.
Для этого достаточно поправить распределение коэффициента \(k_\Gamma\) по высоте проточной части.
При этом потеряется физический смысл коэффициента, но полученное распределение циркуляции лучше соответствует проектному.
Подход, при котором учитывается изменение кольцевой осесимметричной струйки тока и вращение колеса, приводит к более точным результатам профилирования.

\section{Примеры профилирования лопаток различных устройств}

Основное применение разработанного метода заключается в профилировании вентиляторов с меридиональным ускорением потока, но не ограничивается им.
Разбиение течения в проточной части на независимые струйки тока возможно не только для турбомашин с цилиндрической проточной частью, но и со сложной формой меридионального сечения.
Квазитрёхмерное моделирование течения по осесимметричным поверхностям тока в слое переменной толщины методом дискретных вихрей позволяет быстрее сеточных методов определять характеристики нагнетательных турбомашин.

\subsection{Оптимизационные алгоритмы регулируемых осевых вентиляторов}

Совместно с инженерными методами, благодаря относительно высокой скорости выполнения расчёта и низких требований к ЭВМ, разработанный метод профилирования успешно применяется в оптимизационных задачах.
Инженерные методы содержат очень большое количество параметров, что может вызывать появления локальных экстремумов, которые являются проблемой для многих методов математической оптимизации.

В работе \autocite{Zamolodchikov2022} был осуществлен генетический алгоритм для определении формы лопаточных венцов регулируемого вентилятора схемы ВНА+РК+СА.
Суть алгоритма заключается в случайном переборе, комбинировании и вариации параметров \autocite{Holland1992}.
Такой подход позволяет найти глобальный экстремум.

Критерием оптимизации являлось среднее КПД вентилятора \(\eta\) (с учётом веса каждого режима) на заданных режимах работы \(p_v\) и \(Q\).
Начальная популяция состояла из 20 вентиляторов со случайными параметрами.
В результате алгоритмов скрещивания и мутаций получался новый случайны набор комбинаций (популяция).
При этом предпочтение отдаётся вентиляторам с высшей функцией приспособленности (\(\eta\)).
Алгоритм продолжает итеративное выполнение до достижения необходимых параметров, то есть достаточного приближения к глобальному оптимуму.

Для примера работы алгоритма спроектированы два вентилятора с двумя целевыми режимами работы.
Один вентилятор регулировался изменением частоты вращения, об/мин, а другой поворотом лопаток рабочего колеса.
Диаметр РК обоих вентиляторов составлял \(D_\text{к}=0,7\)~м, а относительное удлинение лопатки \(\bar{h}=1\).

Выбранные режимы имели следующие параметры \(p_{\nu 1}=346\)~Па; \(Q_1=2,74\)~м3/с; \(p_{\nu 2}=467\)~Па; \(Q_2 = 5,15\)~м3/с.


\begin{table}[htb]
	\caption{Геометрические параметры двух разработанных регулируемых вентиляторов}
	\label{tab:2vent}
	\centering
	\begin{tabular}{l||c|c|c}
		параметр & варьируемый диапазон & \shortstack{вентилятор 1,\\ \(\theta_\text{РК}=\text{var}\)} & \shortstack{вентилятор 2, \\ \(n=\text{var}\)} \\ \hline \hline
		\(\nu\) & 0,5--0,65 & 0,65 & 0,65 \\ \hline
		\(n\), об/мин & 800--1200 & 1200 & \textbf{1040; 1464} \\ \hline
		\(\bar{f}_\text{РК}\) & 0,02--0,1 & 0,052 & 0,09 \\ \hline
		\(\bar{f}_\text{СА}\) & 0,02--0,1 & 0,052 & 0,0987 \\ \hline
		\(\tau_\text{РК}\) & 0,5--2,0 & 0,545 & 0,55 \\ \hline
		\(\tau_\text{СА}\) & 0,5--2,0 & 1,02 & 0,923 \\ \hline
		\(\theta_\text{РК}\), \textdegree & 10--80 & \textbf{26,7; 45,7} & 33,6 \\ \hline
		\(\theta_\text{СА}\), \textdegree & 10--80 & 78 & 75 \\
	\end{tabular}
\end{table}

Характеристики спроектированных вентиляторов \ref{fig:gleb2022} были получены в проверочном расчёте методом RANS.
Вентилятор, регулируемый поворотом рабочих лопаток, имел КПД на заданных режимах 86\% и 87\%, а регулируемый частотой вращения "--- 82\% и 86\%.

\begin{figure}[htb]
	\centering
	\subcaptionbox{\(p_v\)"---\(Q\)}{\includegraphics[width=0.4\textwidth]{ch4/Gleb2022a}}
	\subcaptionbox{\(\bar{H}\)"---\(\bar{c}_a\)}{\includegraphics[width=0.4\textwidth]{ch4/Gleb2022b}}
	\caption{Характеристики осевых вентиляторов по схеме РК+СА при различных способах регулирования в размерных (а) и безразмерных (б) параметрах.}
	\label{fig:gleb2022}
\end{figure}

\subsection{Направляющий аппарат в S-образном канале}
Высокие требования к топливной экономичности и ограничений по шуму коммерческих самолётов обеспечиваются применением двухконтурных газотурбинных двигателей.
Для обеспечения эффективности таких компоновок используются многовальные системы привода вентилятора и контуров компрессора.
При большой радиальной разнице размеров для их соединения применяют S-образные переходные каналы \autocite{Lu2014}.
Для снижения аэродинамических потерь и обеспечения равномерности на выходе их канала, его конструктивные стойки могут выполняться профилированными и таким образом объединяться со спрямляющим аппаратом.
Эта мера позволяет существенно сократить осевые размеры канала \autocite{An2024}.

Для профилирования спрямляющего аппарата в S-канале можно использовать разработанный метод.
Форма канала была определена ориентируясь на форму испытанного в \autocite{An2024}  канала (рисунок \ref{fig:SShaped}).
Линия совмещения профилей строилась ортогонально линиям тока в меридиональном сечении.
\begin{figure}[htb]
	\centering
	\includegraphics{ch4/SShapedPlane}
	\caption{Меридиональная форма S-канала.}
	\label{fig:SShaped}
\end{figure}

На входе в расчётную область поступает закрученный поток с постоянной по высоте проточной части циркуляцией \(\bar{\Gamma}_1(\bar{h})=\bar{c}_{1u}\bar{r}=0,3\) и коэффициентом расхода \(\bar{c}_{1a}=0,5\).
Расчётное значение циркуляции на выходе из канала \(\Gamma_2(\bar{h})=0\), то есть осевой поток.

Дополнительным требованием стало обеспечение безударного входа по всей высоте лопатки.
Безударный вход означает совпадение передней кромки профиля с первой критической точкой s\textsubscript{1}, в которой поток начинает обтекать лопатку.
Так как все решения полученные методом дискретных вихрей "--- линейно зависимые, поиск безударного режима работы решётки осуществляется по линейной комбинации с коэффициентами \(\lambda\) определёнными по требованию равенства нулю относительной скорости в контрольной точке расположенной на передней кромке профиля \(w_{\text{s}_1}=0\).
Задача сводится к поиску угла атаки и изгиба профиля одновременно удовлетворяющим заданным требованиям.
В таблице приведены параметры полученной решётки и профилей по высоте лопатки.

\begin{table}[htb]
	\caption{Параметры решёток профилей спрямляющего аппарата в S-образном канале.}
	\label{tab:SShapedFoil}
	\begin{center}
		\begin{tabular}{r|c|c|c|c}
			\textnumero~сечения & \(\tau\) & \(i\), \textdegree & \(\bar{f}\) & \(\theta_\text{г}\), \textdegree \\
			\hline
	\hline
	(периферия) 1 & 1.80 & -2.22 & 0.10 & 13.09 \\
	\hline
	2  & 1.85 & -2.10 & 0.10 & 13.00 \\
	\hline
	3  & 1.89 & -2.33 & 0.10 & 13.10 \\
	\hline
	4  & 1.94 & -2.17 & 0.10 & 13.04 \\
	\hline
	(средняя линия) 5  & 1.99 & -2.43 & 0.10 & 13.19 \\
	\hline
	6  & 2.05 & 0.00 & 0.09 & 12.16 \\
	\hline
	7  & 2.11 & 0.00 & 0.09 & 12.23 \\
	\hline
	8  & 2.18 & 0.00 & 0.09 & 12.26 \\
	\hline
	(втулка) 9  & 2.24 & 0.00 & 0.09 & 12.39 \\
		\end{tabular}
	\end{center}
\end{table}

На рисунке \ref{fig:SShapedGamma} изображено сравнение расчётного распределения циркуляции по высоте выхода из канала \(\Gamma_{2\,\text{RANS}}(\bar{h_2})\) на выходе и полученного в проверочном расчёте методом RANS.
\begin{figure}[htb]
	\centering
	\includegraphics{ch4/SShaped_G}
	\caption{Меридиональная форма S-канала}
	\label{fig:SShapedGamma}
\end{figure}
где
\(\bar{a}\) "--- относительное осевое расстояние, \(\bar{a}=a/R_\text{к}\),
\(\bar{p}\) "--- относительное давление на профиле
\begin{equation}
	\bar{p}=\frac{2p}{\rho U^2_\text{к}},
\end{equation}


Относительная средняя интегральная по расходу ошибка величины изменения циркуляции \(\Delta_\Gamma\) составила всего 2\%, где
\begin{equation}
	\bar{\Delta}_\Gamma = \int_{\bar{h}=0}^{\bar{h}=1}\frac{\Gamma_{2\,\text{RANS}}-\Gamma_{2\,\text{Q3D}}}{\Gamma_{2\,\text{RANS}}-\Gamma_1} \, d\bar{c}_m.
		\end{equation}

Обеспечение безударного входа можно оценить по рисунку \ref{fig:SShapedS1}.
В трёх сечениях лопатки \(\bar{h}\) изображено распределение статического давления вблизи передней кромки.
Максимальное давление характеризует расположение на контуре профиля положение точки s\textsubscript{1}.
\begin{figure}[htb]
	\centering
	\subcaptionbox{\(\bar{h}=0,9\)}{\includegraphics{ch4/SShape_s1_1}}
	\subcaptionbox{\(\bar{h}=0,5\)}{\includegraphics{ch4/SShape_s1_2}}
	\subcaptionbox{\(\bar{h}=0,1\)}{\includegraphics{ch4/SShape_s1_3}}
	\caption{Распределение давления у передней кромки лопатки}
	\label{fig:SShapedS1}
\end{figure}

\subsection{Рабочее колесо радиального вентилятора}

\subsection{Рабочее колесо осе-радиального компрессора}

Метод профилирования лопаток рабочих колёс осе-радиальных компрессоров, в отличии от осевых, базируется на статистической обработке экспериментально полученных характеристик \autocite{Becknev1992}.
Важным параметром является коэффициент мощности \(\mu\) "--- параметр характеризующий снижение теоретического напора обусловленного конечным числом лопаток.
В осевых машинах тоже проявляется зависимость от количества лопаток, но принимается что \(\mu=1\).
Его большое значение в радиальных машинах проявляется из-за значительного изменения радиуса поверхности тока и переходе от неинерциальной системе отсчёта в относительном движении к инерционной в абсолютном.

Хорошо развитые моделирования невязкого течения в РК сеточными методами \autocite{Adler1980a} недостаточно подробно описывают характеристики в сравнении с моделированием вязкого \autocite{Adler1980b}.
Однако комбинация методов моделирования в разработанном методе профилирования позволяет относительно быстро получить форму лопатки в сложной меридиональной форме.
Кроме необходимо учитывать переход от моделирования несжимаемого течения к сжимаемому.
Для осевых компрессоров использование характеристик плоских решёток при числах Маха \(\M<0,3\) не приводит к значительной ошибке на расчётном режиме при \(\M < 0,7\) \autocite{Dovjik1968}.

На расчётные параметры (таблица \ref{tab:RadialCompressor}) были спрофилированы лопатки компрессора с проточной частью заданной формы (рисунок \ref{fig:RadialCompressor}), где \(\bar{r}=r/R_\text{к}\) "--- относительный радиус; \(\bar{a}=a/R_\text{к}\) "--- относительная осевая координата.
Относительная толщина профиля в каждом сечении \(\bar{c}=0,02\).
%, коэффициент снижения циркуляции \(k_\Gamma=1.1\).
\begin{table}[htb]
	\caption{Параметры РК осе-радиального компрессора}
	\label{tab:RadialCompressor}
	\begin{center}
		\begin{tabular}{l|c|c}
			Диаметр рабочего колеса, м & \(D_\text{к}\) & 0,265 \\ \hline
			Коэффициент теоретического напора& \(\bar{H}_\text{т}\) &	0,85 \\ \hline
			% Степень повышения давления & \(\pi_\text{к}\) &	2,9 \\ \hline
			Частота вращения, об/мин & \(n\) &	25000 \\ \hline
			Расход воздуха, кг/с & \(G\) &	1,5 \\ \hline
			Коэффициент сужения площади проточной части &\(m_f\)	&	1,22 \\ \hline
			Относительное число Маха на входе & \(\M_{w_1}\)	&	0,42 \\ \hline
			Относительный диаметр периферии на входе& \(\bar{d}_\text{п}\)	&	0,52 \\ \hline
			Относительный диаметр втулки на входе& \(\bar{d}_1\)	&	0,13 \\ \hline
			Закон профилирования &\(c_{u}r\)&\(\const\) \\ \hline
			Количество лопаток &\(Z\)&\(23\)
		\end{tabular}
	\end{center}
\end{table}
\begin{figure}[htb]
	\centering
	\includegraphics{ch4/RadialCompressor}
	\caption{Меридиональная форма осе-радиального компрессора.}
	\label{fig:RadialCompressor}
\end{figure}

Для нормальный условий параметров атмосферы на входе (давление \(p=101325\)~Па, температура \(t=15\)\textcelsius), расчётная степень повышения давления \(\pi^*\), при \(\eta=0,92)\) составляет \(\pi^*=2,68\).

\begin{table}[htb]
	\caption{Параметры решёток профилей осе-радиального компрессора.}
	\label{tab:RadialCompressorFoil}
	\begin{center}
		\begin{tabular}{r|c|c|c|c|c}
			\textnumero~сечения & \(\tau\) & \(i\), \textdegree & \(\bar{f}\) & \(\theta_\text{г}\), \textdegree & \(k_\Gamma\) \\
			\hline \hline
			(периферия) 1 & 4,2 & \(-2\) & 0,1515 & 27,2 &  0,98 \\
			\hline
			(средняя линия) 2 & 5,2 & \(-1,7\) & 0,1457 & 22,8 & 0,98 \\
			\hline
			(втулка) 3 & 8,20 & 2 & \(-0,158\) & 9,1 & 1,03
		\end{tabular}
	\end{center}
\end{table}

Проверочным расчётом выполнено моделирование вязкого сжимаемого течения решением уравнений RANS.
Замыкание выполнялось моделью турбулентности \(k\)"---\(\epsilon\).
На рисунке \ref{fig:RadialCompressorGamma} изображено сравнение расчётного распределения циркуляции по высоте выхода из канала \(\Gamma_{2 \text{RANS}}(\bar{h_2})\) на выходе и полученного в проверочном расчёте RANS.
\begin{figure}[htb]
	\centering
	\includegraphics{ch4/RadialCompressor_G2}
	\caption{Распределение циркуляции потока на выходе из рабочего \(\Gamma_{2}\) колеса по высоте проточной части \(\bar{h}\) осе-радиального компрессора.}
	\label{fig:RadialCompressorGamma}
\end{figure}

%Средняя интегральная по расходу ошибка величины изменения циркуляции \(\Delta_\Gamma\) составила всего 3,5\%.

Более грубые допущения принятые при профилировании лопаток модели несжимаемого течения в отличие от сжимаемого, использованной в RANS-расчёте, приводит к расхождению параметров, которые можно оценить по давлению на профиле (рисунок \ref{fig:RadialCompressor_p}).
Степень повышения давления в компрессоре приводит к значительному повышению плотности воздуха
(рисунок \ref{fig:RadialCompressorRho}), что не учитывается при профилировании.
\begin{figure}[htb]
	\centering
	\subcaptionbox{Давление на профиле \(p\) вдоль осевой координаты \(a\).\label{fig:RadialCompressor_p}}{	\includegraphics[scale=0.9]{ch4/RadialCompressor_p}}
	\subcaptionbox{Плотность воздуха по середине высоты лопатки.\label{fig:RadialCompressorRho}}{	\includegraphics[scale=0.9]{ch4/RadialCompressorRho}}
	
	\caption{Распределения параметров в проточной части.}
\end{figure}

Несмотря на существенное расхождение распределений параметров, интегральные коэффициент теоретического напора \(\bar{H}_\text{т}=0,866\) остался в пределах 1,8\% от расчётного 0,85.
Степень повышения давления составила \(\pi_\text{к}=2,744\).

\section{Заключение}

Разработанная в работе модель течения показала хорошее совпадение распределения давления на лопатке вентилятора с меридиональным ускорением потока с полученными методами решением RANS сеточными методами.
Принимая ограничения накладываемые использованием методов потенциального течения удалось существенно снизить длительность расчёта.
Расчёт распределения давления методом МДВ (600 вихрей) занял 0,2~с.
Время, затраченное на расчёт целого межлопаточного канала RANS (1,4 млн. ячеек), составило 900~с (дольше в 4,5 тысячи раз), а Эйлера 251~c (в 1,26 раз).
Среднее время, затраченное на расчёт вязкого течения в одном сечении (0,1 млн. ячеек), составило 131~с (в 7 раз дольше).
Относительно высокая скорость счёта позволяет отдать преимущество использованию разработанной модели в оптимизационных расчётах.

Модернизированный метод профилирования позволил существенно повысить точность получаемых характеристик, чт является основным применением разработанного аппарата.
В качестве тестовой задачи было проведено профилирование рабочего колеса вентилятора с меридиональным ускорением потока, площадь проточной части которого уменьшалась на 15\% от входа к выходу из РК.
Отклонение от расчётного коэффициента теоретического напора \(\bar{H}_\text{т}\) в тестовой задаче составило 0,2\%.
Это существенно по сравнению с методом профилирования при котором игнорируется изменение формы проточной части, в этом случае отклонение составило 3,8\%.
А профилировании по характеристикам неподвижных решёток пересчитанных на вращающийся венец составила 10\%.

Относительно высокие точность и скорость моделирования позволили успешно задействовать методику профилирования при оптимизации регулируемого осевого вентилятора генетическим алгоритмом.
При регулировании поворотом лопаток рабочего колеса, КПД оптимизированного вентилятора состави 86\% на первом целевом режиме работы, и 87\% на втором.
При регулировании изменением частоты вращения соответственно 82\% и 86\%.

В качестве дополнительного применения рассматривается профилирование лопаток в S-образном канале и осе-радиального компрессора. 

Для заданного на входе в S-канал закрученного потока с \(\bar{\Gamma}_1=0,3\) с коэффициентом расхода \(\bar{c}_{1a}=0,5\) на выходе остаточная закрутка составила всего \(\bar{\Gamma}_2=0,006\), то есть изменение циркуляции \(\bar{\Gamma}_1-\bar{\Gamma}_2\) на 2\% меньше расчётного.
Применение линейно-зависимых результатов моделирования позволило контролировать положение первой критической точке на профиле, и обеспечить безударный вход по большей части лопатки.

Лопатки рабочего колеса осе-центробежного компрессора были спрофилированны по новой методике, что обеспечило совпадение полученных в проверочном моделировании решением RANS характеристик расчётным.
Использование модели несжимаемого газа при профилировании не обеспечило удовлетворительного совпадения эпюр давления на лопатке, но интегральный \(\bar{H}_\text{т}=0,866\) достаточно точно, в пределах 1,8\%, совпал с проверочными.
Степень повышения давления \(\pi_\text{к}=2,744\) превысила ожидаемую на 2,4\%.

\FloatBarrier